\documentclass[11pt]{article}

\usepackage[a4paper,margin=1in]{geometry}
\usepackage{amsmath}
\usepackage{amssymb}
\usepackage{enumitem}
\usepackage{parskip}
\usepackage[hidelinks]{hyperref}

\setlist[itemize]{topsep=2pt,itemsep=1pt,leftmargin=1.4em}
\setlist[enumerate]{topsep=2pt,itemsep=2pt,leftmargin=1.7em}

\title{\vspace{-2em}\textbf{\LARGE Network Position and Academic Performance}\\[0.4em]
\large What We Found So Far}
\author{}
\date{ECC3841 Network Economics --- Presentation 2, Progress Report}

\begin{document}
\maketitle
\vspace{-2.5em}

\section*{Where we left off}

Presentation 1 asked one question: does a student's position in a friendship network
predict their grades, on top of who their direct friends are and on top of raw
popularity? The plan was to measure Katz--Bonacich centrality across a range of the
decay parameter $\phi$, and test it while holding friend-group GPA and popularity
fixed. This brief reports what that test actually found, including a gap in our own
design that we found and fixed along the way.

\section*{What we did}

Four steps, each building on the last.
\begin{enumerate}
  \item Replicate Smirnov and Thurner's two checks, to confirm the data behaves the
        way they reported.
  \item Run centrality on its own, without controlling for popularity: the naive test.
  \item Add popularity as a control and sweep $\phi$, pooling all five student groups:
        the test planned in Presentation 1.
  \item Check that test's assumptions, find one was not actually met, and re-run the
        confirmatory test properly in the one sample where it can be.
\end{enumerate}

\section*{Result 1: the replication holds}

Controlling for a student's own past GPA (coefficient $0.955$, school group,
$n=2{,}088$, $R^2=0.908$), a friend's past GPA does not predict that student's future
GPA: coefficient $-0.0055$, $p=0.557$. New friendships show a smaller GPA gap than
friendships about to end: university, $-0.021$, $p<0.001$; high school, $-0.010$,
$p=0.183$, same direction, not significant. Both match Smirnov and Thurner. This is
the foundation the rest of the brief builds on.

\section*{Result 2: the naive centrality effect is popularity, relabelled}

Katz--Bonacich on its own, net of friend GPA, pooled across all five groups
($n=36{,}696$): coefficient $+0.093$ (standardized), $p<0.001$. Looks like position
matters. Add raw in-degree and out-degree to the same regression: the coefficient
drops to $-0.028$, $p=0.446$, no longer significant, while in-degree itself is
$+0.153$, $p<0.001$. The naive effect was popularity wearing a different name.

\section*{Result 3: the planned headline test, and a gap in it}

Sweeping $\phi$, net of friend GPA and popularity, pooled across all five groups:

\begin{center}
\begin{tabular}{lrrrrrr}
$\phi$ (frac.\ of ceiling) & 0.05 & 0.20 & 0.40 & 0.60 & 0.80 & 0.95 \\
\hline
Katz coefficient & $+0.082$ & $+0.022$ & $-0.034$ & $-0.051$ & $-0.048$ & $-0.036$ \\
$p$-value & $0.012$ & $0.598$ & $0.407$ & $0.115$ & $0.031$ & $0.011$ \\
\end{tabular}
\end{center}

At $\phi=0.85$ (fixed in advance as the headline): coefficient $-0.0415$,
$\mathrm{SE}=0.019$, $p=0.033$, $n=36{,}696$. Read at face value: net of popularity and
friend GPA, position has a small, negative, independent effect.

But this regression never controlled for a student's own past GPA, and it should
have. GPA is time-varying only for the 655 high-school students (6 measurements over
time). For the four university cohorts, \texttt{data.mat} stores a single, static GPA
value per student --- not a column that happens to be constant, but literally one
column. Own past GPA cannot be built as a separate regressor for over 80\% of these
36,696 observations. Presentation 1's model equation promised this control. The test
we actually ran did not have it.

\section*{Result 4: the corrected test}

Restrict the confirmatory test to the one sample where a genuine $t \to t+1$ design
with an own-GPA control is possible: high school, $n=2{,}088$ (535 students, up to 5
transitions each).

Before trusting it, we checked whether adding own past GPA collides with the other
regressors. It does not: it correlates $0.008$ with Katz, $0.065$ with in-degree,
$0.122$ with out-degree, $0.285$ with friend GPA; variance inflation factor $1.11$.
(Katz and in-degree still correlate at $0.866$, VIF $\approx 4$ --- expected, and
exactly what sweeping $\phi$ is for.)

\begin{center}
\begin{tabular}{lrrrrrr}
$\phi$ (frac.\ of ceiling) & 0.05 & 0.20 & 0.40 & 0.60 & 0.80 & 0.95 \\
\hline
Katz coefficient & $-0.004$ & $-0.002$ & $+0.001$ & $+0.003$ & $+0.004$ & $+0.005$ \\
$p$-value & $0.62$ & $0.84$ & $0.93$ & $0.74$ & $0.51$ & $0.36$ \\
\end{tabular}
\end{center}

At every $\phi$, not significant, and the sign is not even consistent. Once a
student's own trajectory is properly controlled, in the one sample where that is
possible, Katz--Bonacich centrality does not predict GPA.

\begin{center}
\fbox{\parbox{0.9\linewidth}{\centering
Net of popularity, friend GPA, and a student's own past performance, we find no
evidence that network position independently predicts grades. The effects that
appeared in Results 2 and 3 do not survive a correctly specified test.}}
\end{center}

\section*{Why this is still worth reporting}

Three specifications, three different pictures: naive centrality looked positive,
the pooled sweep looked negative, the corrected test looks like nothing. The first two
looked like findings and neither survived scrutiny. That progression is itself the
result:
\begin{itemize}
  \item It is the first time this claim --- that Katz--Bonacich centrality
        (Calv\'o-Armengol, Patacchini \& Zenou, 2009) independently predicts
        achievement --- has been tested with a real time lag and an own-performance
        control, which their cross-sectional data could not support.
  \item It shows a general risk: pooling groups with different data structures (GPA
        that varies over time in one, static in four) can manufacture a significant
        coefficient that the one correctly specified subsample does not support.
  \item It is a decision-relevant answer: interventions built around widening
        students' network reach are not supported by this data once the obvious
        confounds --- popularity, friend quality, the student's own trajectory --- are
        properly controlled.
\end{itemize}

\section*{Limitations and what's next}

\begin{itemize}
  \item The corrected test has $n=2{,}088$ from 535 students --- a much wider
        confidence interval than the pooled test's 36,696.
  \item It cannot be extended to the university cohorts: their GPA is a single
        measurement, not a time series, so no lagged design is possible for them with
        this data.
  \item The earlier pattern --- the effect concentrated in sparser networks --- was
        found under the pooled specification, without an own-GPA control. It has not
        been re-checked under the corrected specification and should be treated as
        unconfirmed.
  \item Betweenness and eigenvector centrality were checked in the naive horse race
        only (both similarly explained by degree); neither has been run through the
        corrected, school-only design yet.
\end{itemize}

\vspace{1em}
\noindent\textit{Full study setup, the model, and the identification strategy:
research\_brief.pdf. This brief reports what running that plan actually found.}

\end{document}
